\documentclass[12pt,a4paper]{article}

\usepackage[T1]{fontenc}
\usepackage[utf8]{inputenc}
\usepackage[brazilian]{babel}
\usepackage{geometry}
\usepackage{microtype}
\usepackage{graphicx}
\usepackage{booktabs}
\usepackage{hyperref}
\usepackage{enumitem}

\geometry{margin=2.5cm}
\hypersetup{colorlinks=true,linkcolor=blue,urlcolor=blue,citecolor=blue}

\title{Título do artigo}
\author{Nome do aluno}
\date{2026}

\begin{document}

\maketitle

\begin{abstract}
Apresente o problema, o objetivo, o método e o principal resultado do trabalho.
\end{abstract}

\section{Introdução}
Contextualize o tema, apresente a lacuna e formule a pergunta de pesquisa.

\section{Revisão de literatura}
Organize os trabalhos relacionados e explique como o artigo se posiciona.

\section{Método}
Descreva dados, hipóteses, procedimentos e critérios de análise.

\section{Resultados e discussão}
Apresente os resultados e discuta suas implicações e limitações.

\section{Conclusão}
Retome a pergunta de pesquisa e indique os próximos passos.

\begin{thebibliography}{9}
\bibitem{exemplo}
AUTOR. \emph{Título da referência}. Editora ou periódico, ano.
\end{thebibliography}

\end{document}
